\documentclass[14pt]{beamer}
\usepackage{pepslides}
\title{Waveform tomography}
\lecturenumber{20}
\hypersetup{pdftitle={Lecture 20: Waveform tomography}}


\begin{document}

\titleslide

\begin{frame}[t]
\heading{From last time\ldots}
Using Fermat's principle, we find that:
\begin{eqn}
\delta T=-\int_{\mathrm{ray}}\frac{\delta v}{v^{2}}\dd s.
\end{eqn}
Expand the velocity in a finite-dimensional basis:
\begin{eqn}
\delta\alpha(\mathbf{x})=\sum_{j=1}^{m}\delta\alpha_{j}\varphi_{j}(\mathbf{x}).
\end{eqn}
Formulation as a regularised least squares problem:
\begin{eqn}
J(\mathbf{m})=\frac{1}{2}(\mathbf{A}\mathbf{m}-\mathbf{d})^{T}\mathbf{C}^{-1}(\mathbf{A}\mathbf{m}-\mathbf{d})+\frac{\lambda}{2}\mathbf{m}^{T}\mathbf{B}\mathbf{m}.
\end{eqn}
\end{frame}

\begin{frame}
\centering
\animation[width=0.85\textwidth]{healing_t1.pdf}{healing.webm}

Wavefront healing as a plane wave passes a small heterogeneity.

Scattering regime: {\eqsize$\dfrac{L}{\lambda}\approx1$}\hspace{4em}
Geometric optics regime: {\eqsize$\dfrac{L}{\lambda}\gg1$}
\end{frame}

\begin{frame}
\centering
\animation[width=0.85\textwidth]{healing_t2.pdf}{healing.webm}

\phantom{Wavefront healing as a plane wave passes a small heterogeneity.}

Scattering regime: {\eqsize$\dfrac{L}{\lambda}\approx1$}\hspace{4em}
Geometric optics regime: {\eqsize$\dfrac{L}{\lambda}\gg1$}
\end{frame}

\begin{frame}
\centering
\animation[width=0.85\textwidth]{healing_t3.pdf}{healing.webm}

\phantom{Wavefront healing as a plane wave passes a small heterogeneity.}

Scattering regime: {\eqsize$\dfrac{L}{\lambda}\approx1$}\hspace{4em}
Geometric optics regime: {\eqsize$\dfrac{L}{\lambda}\gg1$}
\end{frame}

\begin{frame}
Typical periods for global seismology: 1--40\,s.

Corresponding wavelengths: 10--400\,km.

\centering
\includegraphics[height=0.55\paperheight]{farnetani1997_plumes.png}
\credit{Farnetani (1997)}
\end{frame}

\sectionslide{Waveform inversion}

\begin{frame}
Equations of motion:
\begin{eqn}
\begin{gathered}
\rho\frac{\partial^{2}u_{i}}{\partial t^{2}}-\frac{\partial}{\partial x_{j}}\left(A_{ijkl}\frac{\partial u_{k}}{\partial x_{l}}-\overline{S}_{ij}\right)=0,\\[\eqgap]
\left(A_{ijkl}\frac{\partial u_{k}}{\partial x_{l}}-\overline{S}_{ij}\right)\hat{n}_{j}=0.
\end{gathered}
\end{eqn}
These can be solved numerically in realistic laterally heterogeneous earth models using
techniques like the \textbf{finite element method}.

These are demanding calculations, run on HPC systems using parallel computing over
hundreds or thousands of cores. GPU computing is now common.

Each simulation takes \textbf{tens of minutes to hours} depending on resources.
\end{frame}

\begin{frame}
\centering
\includegraphics[height=0.7\paperheight]{L8F1.pdf}

\raggedright
A typical data set: around $10^{3}$ earthquakes recorded at thousands of stations, each earthquake
requiring its own simulation.
\credit{Cui \emph{et al.} (2024)}
\end{frame}

\begin{frame}
\centering
\includegraphics[width=0.8\textwidth]{L8F2.png}
\begin{eqn}
J=\frac{1}{2}\int_{0}^{T}\|\mathbf{u}(\mathbf{x}_{r},t)-\mathbf{u}^{\mathrm{obs}}(t)\|^{2}\dd t.
\end{eqn}
\credit{Tape \emph{et al.} (2010)}
\end{frame}

\begin{frame}
Misfit is an implicit function of the parameters in the elastic wave equation:
\begin{enumerate}
\item \textbf{Density} -- piecewise smooth function
\item \textbf{Elastic tensor} -- piecewise smooth tensor field (21 independent components)
\item \textbf{Stress glut} -- space and time dependent tensor field
\item \textbf{Geometry} of the model (e.g.\ topography on internal boundaries)
\end{enumerate}
Relation can be condensed to:
\begin{eqn}
\mathbf{m}\mapsto J(\mathbf{m}).
\end{eqn}
The aim is to \textbf{minimise} this functional.

Strictly defined on an \textbf{infinite-dimensional space}. In practice reduced to a finite but
high-dimensional space linked to the numerical discretisation of the equations of motion.

\textbf{Millions or billions of degrees of freedom would be typical.}
\end{frame}

\begin{frame}
\begin{columns}[c]
\begin{column}{0.5\textwidth}
\includegraphics[width=\linewidth]{gradient_descent_contours.png}
\end{column}
\begin{column}{0.44\textwidth}
Iterative reduction of the misfit using the local gradient to take steps downhill.

For large-scale optimisation problems, this is often the only viable approach.

Bayesian sampling methods are not typically an option, though this is an active research topic.

To proceed, we need a way of evaluating the gradients.
\end{column}
\end{columns}
\end{frame}

\begin{frame}
Functional derivative defined through a first-order Taylor expansion:
\begin{eqn}
J(\mathbf{m}+\delta\mathbf{m})=J(\mathbf{m})+\langle DJ(\mathbf{m})|\delta\mathbf{m}\rangle+O(\|\delta\mathbf{m}\|^{2}).
\end{eqn}
The terms \textbf{gradient} and \textbf{sensitivity kernel} are equivalent and widely used.

Focusing on S-wave velocity variations alone, we can write to first order:
\begin{eqn}
\langle DJ(\mathbf{m})|\delta\mathbf{m}\rangle=\int_{M}K_{\beta}\,\delta\beta\dd^{3}\mathbf{x}.
\end{eqn}
Methods for calculation will be discussed later.
\end{frame}

\begin{frame}
Condition for a (local) minimum, and the condition that follows from it:
\begin{eqn}
\begin{gathered}
J(\overline{\mathbf{m}}+\delta\mathbf{m})\ge J(\overline{\mathbf{m}}),\\[\eqgap]
DJ(\overline{\mathbf{m}})=\mathbf{0}.
\end{gathered}
\end{eqn}
Starting at an initial point, we consider the update:
\begin{eqn}
\mathbf{m}_{1}=\mathbf{m}_{0}-\lambda\,DJ(\mathbf{m}_{0}).
\end{eqn}
Step length determined by approximately solving a 1D minimisation problem:
\begin{eqn}
f(\lambda)=J[\mathbf{m}_{0}-\lambda\,DJ(\mathbf{m}_{0})].
\end{eqn}
Process is repeated until convergence or the misfit is sufficiently small. This is the \textbf{method of steepest descent}.
\end{frame}

\begin{frame}
\centering
\includegraphics[height=0.55\paperheight]{local_minima_surface.png}

\raggedright
Gradient-based methods produce only \textbf{local minima}.

Uncertainty quantification possible to a limited extent, but adds significantly to the cost.
\end{frame}

\sectionslide{Calculating the gradients}

\begin{frame}
\textbf{Finite-difference approximation:}
\begin{eqn}
\langle DJ(\mathbf{m})|\delta\mathbf{m}\rangle\approx J(\mathbf{m}+\delta\mathbf{m})-J(\mathbf{m}).
\end{eqn}
Approximation of the gradient relative to an $m$-dimensional \textbf{model parameterisation} requires
$m+1$ \textbf{forward calculations}.

One gradient calculation using this method for a realistic problem would take of order
hundreds of years on a massive HPC system!

The \textbf{adjoint method} provides the exact gradient at a cost of \textbf{two forward simulations}.
\end{frame}

\begin{frame}
We want to minimise:
\begin{eqn}
J=\frac{1}{2}\int_{0}^{T}\|\mathbf{u}(\mathbf{x}_{r},t)-\mathbf{u}^{\mathrm{obs}}(t)\|^{2}\dd t,
\end{eqn}
subject to the constraints:
\begin{eqn}
\begin{gathered}
\rho\frac{\partial^{2}u_{i}}{\partial t^{2}}-\frac{\partial}{\partial x_{j}}\left(A_{ijkl}\frac{\partial u_{k}}{\partial x_{l}}-\overline{S}_{ij}\right)=0,\\[\eqgap]
\left(A_{ijkl}\frac{\partial u_{k}}{\partial x_{l}}-\overline{S}_{ij}\right)\hat{n}_{j}=0.
\end{gathered}
\end{eqn}
This can be done using the \textbf{method of Lagrange multipliers}.
\end{frame}

\begin{frame}
We start by defining the \textbf{Lagrangian}:
\begin{eqnsmall}
\begin{aligned}
L(\mathbf{u},\mathbf{u}',\mathbf{w}',\mathbf{m}) = J(\mathbf{u}) &- \int_{0}^{T}\int_{M}\left[\rho\frac{\partial^{2}u_{i}}{\partial t^{2}}-\frac{\partial}{\partial x_{j}}\left(A_{ijkl}\frac{\partial u_{k}}{\partial x_{l}}-\overline{S}_{ij}\right)\right]u_{i}'\dd^{3}\mathbf{x}\dd t\\[\eqgap]
&- \int_{0}^{T}\int_{\partial M}\left(A_{ijkl}\frac{\partial u_{k}}{\partial x_{l}}-\overline{S}_{ij}\right)\hat{n}_{j}w_{i}'\dd S\dd t.
\end{aligned}
\end{eqnsmall}
Each constraint is associated with a \textbf{Lagrange multiplier field}.

Zero initial conditions on displacement and velocity are being imposed directly.

Within the Lagrangian, all fields are regarded as being independent.
\end{frame}

\begin{frame}
The \textbf{Lagrange multiplier theorem} states that the following equality holds:
\begin{eqn}
DJ(\mathbf{m})=D_{\mathbf{m}}L(\mathbf{u},\mathbf{u}',\mathbf{w}',\mathbf{m}),
\end{eqn}
so long as the following conditions are met:
\begin{eqn}
\begin{aligned}
D_{\mathbf{u}}L(\mathbf{u},\mathbf{u}',\mathbf{w}',\mathbf{m})&=0,\\
D_{\mathbf{u}'}L(\mathbf{u},\mathbf{u}',\mathbf{w}',\mathbf{m})&=0,\\
D_{\mathbf{w}'}L(\mathbf{u},\mathbf{u}',\mathbf{w}',\mathbf{m})&=0.
\end{aligned}
\end{eqn}
Subscripts used for partial functional derivatives.

You only need to know how to apply this method, \textbf{not to prove} that it works -- details are
provided in an appendix to the handout for anyone interested.
\end{frame}

\begin{frame}
Recall that the Lagrangian is:
\begin{eqnsmall}
\begin{aligned}
L(\mathbf{u},\mathbf{u}',\mathbf{w}',\mathbf{m}) = J(\mathbf{u}) &- \int_{0}^{T}\int_{M}\left[\rho\frac{\partial^{2}u_{i}}{\partial t^{2}}-\frac{\partial}{\partial x_{j}}\left(A_{ijkl}\frac{\partial u_{k}}{\partial x_{l}}-\overline{S}_{ij}\right)\right]u_{i}'\dd^{3}\mathbf{x}\dd t\\[\eqgap]
&- \int_{0}^{T}\int_{\partial M}\left(A_{ijkl}\frac{\partial u_{k}}{\partial x_{l}}-\overline{S}_{ij}\right)\hat{n}_{j}w_{i}'\dd S\dd t.
\end{aligned}
\end{eqnsmall}
Stationarity with respect to $\mathbf{u}'$ gives the equations of motion.

Stationarity with respect to $\mathbf{w}'$ gives the boundary conditions.

This is how the \textbf{forward problem} emerges within the method.
\end{frame}

\begin{frame}
Varying with respect to $\mathbf{u}$ gives:
\begin{eqnsmall}
\begin{aligned}
\langle D_{\mathbf{u}}L(\mathbf{u},\mathbf{u}',\mathbf{w}',\mathbf{m})|\delta\mathbf{u}\rangle &= \langle DJ(\mathbf{u})|\delta\mathbf{u}\rangle\\
&\quad - \int_{0}^{T}\int_{M}\left[\rho\frac{\partial^{2}\delta u_{i}}{\partial t^{2}}-\frac{\partial}{\partial x_{j}}\left(A_{ijkl}\frac{\partial\delta u_{k}}{\partial x_{l}}\right)\right]u_{i}'\dd^{3}\mathbf{x}\dd t\\
&\quad - \int_{0}^{T}\int_{\partial M}\left(A_{ijkl}\frac{\partial\delta u_{k}}{\partial x_{l}}\right)\hat{n}_{j}w_{i}'\dd S\dd t.
\end{aligned}
\end{eqnsmall}
From the chosen form for $J$ we can write:
\begin{eqn}
\begin{gathered}
\langle DJ(\mathbf{u})|\delta\mathbf{u}\rangle=\int_{0}^{T}\int_{\partial M}h_{i}'\delta u_{i}\dd S\dd t,\\[\eqgap]
h_{i}'(\mathbf{x},t)=[u_{i}(\mathbf{x}_{r},t)-u_{i}^{\mathrm{obs}}(t)]\delta(\mathbf{x}-\mathbf{x}_{r}).
\end{gathered}
\end{eqn}
\end{frame}

\begin{frame}
We need the following condition to hold:
\begin{eqnsmall}
\begin{aligned}
\int_{0}^{T}\int_{\partial M}h_{i}'\delta u_{i}\dd S\dd t &- \int_{0}^{T}\int_{M}\left[\rho\frac{\partial^{2}\delta u_{i}}{\partial t^{2}}-\frac{\partial}{\partial x_{j}}\left(A_{ijkl}\frac{\partial\delta u_{k}}{\partial x_{l}}\right)\right]u_{i}'\dd^{3}\mathbf{x}\dd t\\[\eqgap]
&- \int_{0}^{T}\int_{\partial M}\left(A_{ijkl}\frac{\partial\delta u_{k}}{\partial x_{l}}\right)\hat{n}_{j}w_{i}'\dd S\dd t=0,
\end{aligned}
\end{eqnsmall}
for all variations of the displacement vector.

To see the consequence, we need to pass all the derivatives to the other terms using
integration by parts in time and the divergence theorem spatially.

Full details are in the handout.
\end{frame}

\begin{frame}
The result of the long calculation is that:
\begin{eqnsmall}
\begin{aligned}
-&\int_{0}^{T}\int_{M}\left[\rho\frac{\partial^{2}u_{i}'}{\partial t^{2}}-\frac{\partial}{\partial x_{j}}\left(A_{ijkl}\frac{\partial u_{k}'}{\partial x_{l}}\right)\right]\delta u_{i}\dd^{3}\mathbf{x}\dd t\\
-&\int_{0}^{T}\int_{\partial M}\left(A_{ijkl}\frac{\partial\delta u_{k}}{\partial x_{l}}\right)\hat{n}_{j}(w_{i}'-u_{i}')\dd S\dd t\\
+&\int_{0}^{T}\int_{\partial M}\left[h_{i}'-\left(A_{ijkl}\frac{\partial u_{k}'}{\partial x_{l}}\right)\hat{n}_{j}\right]\delta u_{i}\dd S\dd t\\
-&\int_{M}\rho\left\{\left[\frac{\partial\delta u_{i}}{\partial t}u_{i}'\right]_{t=T}-\left[\delta u_{i}\frac{\partial u_{i}'}{\partial t}\right]_{t=T}\right\}\dd^{3}\mathbf{x}=0.
\end{aligned}
\end{eqnsmall}
From this we can read off the necessary adjoint equations.
\end{frame}

\begin{frame}
We have shown that the Lagrange multipliers must satisfy the equations:
\begin{eqn}
\begin{gathered}
\rho\frac{\partial^{2}u_{i}'}{\partial t^{2}}-\frac{\partial}{\partial x_{j}}\left(A_{ijkl}\frac{\partial u_{k}'}{\partial x_{l}}\right)=0,\\[\eqgap]
\left(A_{ijkl}\frac{\partial u_{k}'}{\partial x_{l}}\right)\hat{n}_{j}=h_{i}', \qquad
u_{i}'(\mathbf{x},T)=\frac{\partial u_{i}'}{\partial t}(\mathbf{x},T)=0,
\end{gathered}
\end{eqn}
along with the equality $\mathbf{w}'=\mathbf{u}'$ on the boundary.

This is a \textbf{terminal-value problem} for $\mathbf{u}'$ driven by the boundary traction $h_{i}'$:
\begin{eqn}
h_{i}'(\mathbf{x},t)=[u_{i}(\mathbf{x}_{r},t)-u_{i}^{\mathrm{obs}}(t)]\delta(\mathbf{x}-\mathbf{x}_{r}).
\end{eqn}
\end{frame}

\begin{frame}
The terminal-value problem must be solved backwards in time from $t=T$. Instead, define the time-reversed field:
\begin{eqn}
u_{i}^{\dagger}(\mathbf{x},t)=u_{i}'(\mathbf{x},T-t).
\end{eqn}
This satisfies the same wave equation, with boundary traction and initial conditions:
\begin{eqn}
\begin{gathered}
\left(A_{ijkl}\frac{\partial u_{k}^{\dagger}}{\partial x_{l}}\right)\hat{n}_{j}=h_{i}^{\dagger}, \quad
h_{i}^{\dagger}(\mathbf{x},t)=h_{i}'(\mathbf{x},T-t),\\[\eqgap]
u_{i}^{\dagger}(\mathbf{x},0)=\frac{\partial u_{i}^{\dagger}}{\partial t}(\mathbf{x},0)=0.
\end{gathered}
\end{eqn}
An \textbf{initial-value problem} of the same form as the forward problem, solved with the same numerical code.

Time reversal is just a device for converting a terminal-value problem into an initial-value one.
\end{frame}

\begin{frame}
Using the multiplier theorem and integrating by parts, we find that:
\begin{eqnsmall}
\langle DJ(\mathbf{m})|\delta\mathbf{m}\rangle=\int_{0}^{T}\int_{M}\left[\delta\rho\frac{\partial u_{i}}{\partial t}\frac{\partial u_{i}'}{\partial t}-\delta A_{ijkl}\frac{\partial u_{i}'}{\partial x_{j}}\frac{\partial u_{k}}{\partial x_{l}}+\delta\overline{S}_{ij}\frac{\partial u_{i}'}{\partial x_{j}}\right]\dd^{3}\mathbf{x}\dd t.
\end{eqnsmall}
We can then read off the \textbf{sensitivity kernels}:
\begin{eqn}
\begin{aligned}
K_{\rho}(\mathbf{x})&=\int_{0}^{T}\frac{\partial u_{i}}{\partial t}(\mathbf{x},t)\frac{\partial u_{i}'}{\partial t}(\mathbf{x},t)\dd t,\\[\eqgap]
K_{A_{ijkl}}(\mathbf{x})&=-\int_{0}^{T}\frac{\partial u_{i}'}{\partial x_{j}}(\mathbf{x},t)\frac{\partial u_{k}}{\partial x_{l}}(\mathbf{x},t)\dd t,\\[\eqgap]
K_{\overline{S}_{ij}}(\mathbf{x},t)&=\frac{\partial u_{i}'}{\partial x_{j}}(\mathbf{x},t).
\end{aligned}
\end{eqn}
\end{frame}

\begin{frame}
In terms of the time-reversed field, the kernels read:
\begin{eqn}
\begin{aligned}
K_{\rho}(\mathbf{x})&=-\int_{0}^{T}\frac{\partial u_{i}}{\partial t}(\mathbf{x},t)\frac{\partial u_{i}^{\dagger}}{\partial t}(\mathbf{x},T-t)\dd t,\\[\eqgap]
K_{A_{ijkl}}(\mathbf{x})&=-\int_{0}^{T}\frac{\partial u_{i}^{\dagger}}{\partial x_{j}}(\mathbf{x},T-t)\frac{\partial u_{k}}{\partial x_{l}}(\mathbf{x},t)\dd t,\\[\eqgap]
K_{\overline{S}_{ij}}(\mathbf{x},t)&=\frac{\partial u_{i}^{\dagger}}{\partial x_{j}}(\mathbf{x},T-t).
\end{aligned}
\end{eqn}
The forward field at time $t$ is paired with the adjoint field after it has evolved for a time $T-t$.

Here $\mathbf{u}^{\dagger}$ is the \textbf{adjoint wavefield} and $h_{i}^{\dagger}$ the \textbf{adjoint traction} (or \textbf{adjoint source}).
\end{frame}

\begin{frame}
For each gradient calculation we must:
\begin{enumerate}
\item Solve the forward problem for $\mathbf{u}$
\item From this solution, compute $J$ and form the adjoint traction $h_{i}^{\dagger}$
\item Solve the adjoint problem for $\mathbf{u}^{\dagger}$ as an initial-value problem
\item Combine the forward and adjoint fields to form the sensitivity kernels
\end{enumerate}
Such a calculation has the cost of two forward calculations (three in practice, so that the
forward field need not be stored).

The adjoint equations are identical in form to the forward ones, and so the same numerical code can be used.

Choice of misfit enters into the problem only through the form of the adjoint traction.

\textbf{Calculations that would take hundreds of years are reduced to minutes or hours.}
\end{frame}

\begin{frame}
In practice, delay times are measured using \textbf{cross-correlations} of windowed waveforms:
\begin{eqn}
\begin{gathered}
s^{\mathrm{obs}}(t)=w(t)\hat{\nu}_{i}u_{i}^{\mathrm{obs}}(t), \qquad s(\mathbf{x}_{r},t)=w(t)\hat{\nu}_{i}u_{i}(\mathbf{x}_{r},t),\\[\eqgap]
C(\tau)=\int_{-\infty}^{\infty}s^{\mathrm{obs}}(t-\tau)s(\mathbf{x}_{r},t)\dd t.
\end{gathered}
\end{eqn}
The time shift $\overline{\tau}$ that maximises the cross-correlation is taken as the delay time.

A late observed arrival gives $\overline{\tau}<0$, the opposite sign to the delay time of Lecture~19.
Only the perturbation $\delta\overline{\tau}$ enters the kernels, so nothing depends on this choice.
\end{frame}

\begin{frame}
This defines a functional of the synthetic waveform, and hence we can derive the
associated sensitivity kernels using the adjoint method.

Perturbing the condition $\ddns C/\ddns\tau(\overline{\tau})=0$, with the waveforms differing only by the time shift:
\begin{eqn}
\delta\overline{\tau}=-\frac{1}{N}\int_{-\infty}^{\infty}\dot{s}(\mathbf{x}_{r},t)\,\delta s(\mathbf{x}_{r},t)\dd t,
\quad N=\int_{-\infty}^{\infty}\dot{s}(\mathbf{x}_{r},t)^{2}\dd t.
\end{eqn}
From this we read off the adjoint traction:
\begin{eqn}
h_{i}'(\mathbf{x},t)=-\frac{1}{N}w(t)\dot{s}(\mathbf{x}_{r},t)\hat{\nu}_{i}\delta(\mathbf{x}-\mathbf{x}_{r}).
\end{eqn}
The observed seismogram has dropped out: the kernel depends only on the reference model.
\end{frame}

\begin{frame}
For the cross-correlation delay time, we obtain to first order:
\begin{eqn}
\delta\overline{\tau}=\int_{M}(K_{\alpha}\,\delta\alpha+K_{\beta}\,\delta\beta)\dd^{3}\mathbf{x}.
\end{eqn}
This can be compared with the ray-theoretic result from Fermat's principle:
\begin{eqn}
\delta T=-\int_{\mathrm{ray}}\frac{\delta c}{c^{2}}\dd s.
\end{eqn}
Within ray theory all the sensitivity lies on the ray path. What does the adjoint method give?
\end{frame}

\begin{frame}
Building a P-wave kernel by the adjoint method in two dimensions:
\begin{eqn}
K_{\alpha}(\mathbf{x})=-2\rho\alpha\int_{0}^{T}\nabla\cdot\mathbf{u}(\mathbf{x},t)\,
\nabla\cdot\mathbf{u}^{\dagger}(\mathbf{x},T-t)\dd t.
\end{eqn}
\centering
\animation[width=0.92\textwidth]{L8F3.pdf}{L8F3.webm}

\raggedright
Sensitivity accumulates only where the forward field at time $t$ overlaps the adjoint field at time $T-t$.
\end{frame}

\begin{frame}
\centering
\includegraphics[width=\textwidth]{L8F4.pdf}

\raggedright
The kernel is negative throughout the \textbf{first Fresnel zone}, since an increase in wave speed there
advances the arrival, and changes sign in the surrounding zones.
\end{frame}

\begin{frame}
\centering
\includegraphics[width=\textwidth]{taylor_test.pdf}

\raggedright
A \textbf{Taylor test}: a fixed perturbation $\delta\alpha$ is scaled by $s$, and the change in the delay time found by
rerunning the forward problem is compared with the prediction from the kernel. They agree to first order in $s$, the
difference shrinking like $s^{2}$.
\end{frame}

\begin{frame}
\centering
\includegraphics[height=0.68\paperheight]{bdkernels_p.pdf}

\raggedright
In three dimensions the sensitivity lies within the first Fresnel zone, which narrows at shorter periods, but
curiously vanishes on the ray path itself. P-wave traveltime at $60^{\circ}$ in PREM, measured in the window shaded.
\credit{Calculated with YSpec}
\end{frame}

\begin{frame}
\centering
\includegraphics[width=0.95\textwidth]{bdkernels_s.pdf}

\raggedright
Named \textbf{banana doughnut kernels} by Tony Dahlen (1942--2007). S-wave traveltime at $60^{\circ}$ in PREM.
\credit{Calculated with YSpec}
\end{frame}

\begin{frame}
\centering
\includegraphics[width=0.95\textwidth]{bdkernels_ss.pdf}

\raggedright
The kernel of SS, reflected at the surface midway between the source and the receiver, is far from a single banana.
\credit{Calculated with YSpec}
\end{frame}

\begin{frame}
\begin{columns}[c]
\begin{column}{0.48\textwidth}
\centering
\includegraphics[height=0.9\paperheight]{L8F6.pdf}
\end{column}
\begin{column}{0.46\textwidth}
GLAD-M35: shear wave speed, bulk sound speed and compressional wave speed at depths from 140 to 2800\,km.

Waveforms from 2160 earthquakes at 11\,575 stations, and 35 iterations of gradient-based optimisation.
\end{column}
\end{columns}
\credit{Cui \emph{et al.} (2024)}
\end{frame}

\begin{frame}
\begin{columns}[c]
\begin{column}{0.48\textwidth}
\centering
\includegraphics[height=0.9\paperheight]{L8F7.pdf}
\end{column}
\begin{column}{0.46\textwidth}
Uncertainty can be estimated approximately: the Hessian of the misfit plays the role of the inverse posterior covariance.

Its action on a vector costs a few further forward and adjoint simulations.

Uncertainty is reduced most where the data coverage is best.
\end{column}
\end{columns}
\credit{Cui \emph{et al.} (2024)}
\end{frame}

\begin{frame}
\begin{columns}[c]
\begin{column}{0.56\textwidth}
\centering
\includegraphics[height=0.9\paperheight]{L8F8.pdf}
\end{column}
\begin{column}{0.4\textwidth}
Six global models at 2800\,km depth, three of them from waveform tomography.

They agree on the largest features, but differ considerably in detail, notably in the bulk sound speed.
\end{column}
\end{columns}
\credit{Cui \emph{et al.} (2024)}
\end{frame}

\begin{frame}
\centering
\includegraphics[height=0.7\paperheight]{L8F9.pdf}

\raggedright
Cross-sections beneath hot spots: the models differ markedly in whether they show continuous
low-velocity conduits, or plumes, rising through the mantle.
\credit{Cui \emph{et al.} (2024)}
\end{frame}

\begin{frame}[t]
\heading{What you need to know and be able to do}
\begin{enumerate}
\item[(i)] That waveform tomography depends on the ability to efficiently solve the elastic wave equation in realistic
earth models.
\item[(ii)] How least-squares waveform misfits can be defined, and what is meant by their gradients.
\item[(iii)] How gradient-based optimisation works, why it is useful, and what its limitations are.
\item[(iv)] How the adjoint method can be used to calculate gradients, and what its advantages are. The derivation in
this lecture is too lengthy for you to reproduce in full. But you may be asked to perform parts of the calculations
with suitable guidance or to apply the method to simpler systems.
\end{enumerate}
\end{frame}

\end{document}
